\subsection{The actual small local image of the mixed braid evaluation}

\begin{result}{A.4.27}{Proved}{The local signed mixed target is an explicit Heisenberg extension of order 432}
Use the literal local $B_4$ support in A.4.26, embedded in
$\Pi_{1,9}$ by retaining one named puncture and the handle and
fixing the other seven puncture generators.  Restrict its actual
action to the faithful thirty-six-dimensional quaternion homology
of A.4.15, and join the ordinary compact period modulo four.
Represent the latter by its permutation of the sixteen vectors
of $(\mathbb Z/4)^2$.  This gives the explicit finite homomorphism
\[
 q_{\rm loc}:B_4\longrightarrow
            \operatorname{GL}_{52}(\mathbb F_3),
 \qquad \sigma_i\longmapsto\operatorname{diag}(M_i,P_i),
                                                               \tag{LI.1}
\]
where $M_i$ and $P_i$ are retained matrices obtained from the
literal free-group substitutions below.  It has image order
$93312$.  Put $D=q_{\rm loc}(d)$ and
$V=q_{\rm loc}((aba)^2)$.  Their generated subgroup $Q$ has
order $432$, with
\[
 [Q,Q]=N\cong\operatorname{UT}_3(\mathbb F_3),\qquad
 Q/N\cong C_4^2,\qquad Z(Q)\cong C_{12}.                \tag{LI.2}
\]
The element $z=(DV)^2$ is central of order six.  Both $D$ and
$V$ have order four, and $DV$ has order twelve.

Here are exact coordinates, including an actual splitting.
Write an element of $N$ as $(i,j,k)\in\mathbb F_3^3$, with
\[
 (i,j,k)(i',j',k')=(i+i',j+j',k+k'-ji').                 \tag{LI.3}
\]
There are commuting actual elements $u,v$ of order four,
with $u=D$, whose images form the stated $C_4^2$ quotient.
In the normal form $(i,j,k)u^p v^q$, the original two
mixed generators and their central product are
\[
\begin{array}{c|c}
 D&(0,0,0;1,0)\\
 V&(2,2,2;0,1)\\
 (DV)^2&(0,0,2;2,2).
\end{array}                                                    \tag{LI.4}
\]
Both $u$ and $v$ act on the three Heisenberg generators
$x=(1,0,0)$, $y=(0,1,0)$, $z_N=(0,0,1)$ by
\[
 x\longmapsto(2,1,2),\qquad
 y\longmapsto(1,1,2),\qquad z_N\longmapsto z_N.         \tag{LI.5}
\]
The two actual free-group evaluations required in (ABR.4) are
\[
\begin{array}{c|cc}
 &(X\mapsto)&(Y\mapsto)\\ \hline
 f(D,V^2)&(0,0,0;1,0)&(2,0,2;0,2)\\
 f(D^2,V)&(0,0,0;2,0)&(2,2,2;0,1).
\end{array}                                                    \tag{LI.6}
\]
All $432$ group elements, their actual matrix words, all generator
transitions, the $729$ Heisenberg multiplication rules, and the
splitting in (LI.4)--(LI.5) are independently enumerated in the
certificate.
For a commutator proword satisfying antisymmetry, write its two
values in (LI.6) as $(i,j,k_1)$ and $(i,j,k_2)$ in $N$.
The first two coordinates necessarily agree.  The entire local
weight test (ABR.4) is precisely the remaining scalar equation
\[
                         k_1-k_2=2c\quad\hbox{in }\mathbb F_3.
                                                               \tag{LI.9}
\]
  No coordinate tuple is inferred to lift to a
Grothendieck--Teichm\"uller element.
\end{result}

\begin{proof}
The four-involution Hurwitz model used for the Ferrari
pointpush calculation has even kernel with basis
\[
 A=v_0,\quad B=u_0^{-1},\quad
 C=u_0^{-1}v_0u_0v_0^{-1}w_0v_0^{-1}u_0,
\]
where $u_0=x_2x_1$, $v_0=x_3x_1$, $w_0=x_4x_1$.
Each adjacent Hurwitz substitution
$x_i\mapsto x_ix_{i+1}x_i$, $x_{i+1}\mapsto x_i$
induces a literal automorphism on this basis.  It fixes
$[A,B]C$ and preserves the conjugacy class of $C$.
Extending it by the identity on $C_2,\ldots,C_8$ therefore
gives the actual local surface automorphism used in (LI.1).
The source program verifies these free-word identities and
stores all three complete ten-generator substitutions.

Their images under the retained quaternion map give all
$24$ automorphisms of $Q_8$.  The complete $72$ Schreier
words of this residue action restrict to the same primitive
multiplicity space $M_{18}$ as A.4.15.  Their common fixed
space has dimension fourteen, and their difference images
span a four-dimensional subspace $U$.  The restriction of
the actual symmetric form to $U$ has rank two.  Thus its
radical has dimension two.  The actual quaternion-kernel
matrix group has order $972=4\cdot3^5$.
The finite matrix computation uses the original $18$-dimensional
matrices; its GAP input and output are retained, so this
count does not replace the source action by an abstract
orthogonal group.

For the signed target the program instead retains the full
$36$-dimensional faithful module, so the nontrivial quaternion
residues remain visible.  Its ordinary two-dimensional period
matrix modulo four acts as a permutation of sixteen actual
vectors, giving the second block in (LI.1).  Exact finite group
calculation on these three $52$-dimensional matrices gives the
full image order $93312$ and its solvability.  Restricting to
the actual words $d$ and $(aba)^2$ gives order $432$.
The program obtains all these orders independently from the
same literal substitutions; the smaller mixed enumeration
below does not require GAP.

\par\medskip\noindent\textbf{Complete independent enumeration of the mixed group.}\quad
Starting with the identity matrix and multiplying on the right
by $D,V$ produces exactly $432$ distinct matrices.  The program
records a word for each one and both outgoing generator edges.
The two exponent sums of each word modulo four agree on every
edge, giving an actual surjection $Q\to C_4^2$ with kernel of
order twenty-seven.  All its elements have order dividing three.
It contains two explicit noncommuting elements $x,y$; setting
$z_N=[x,y]$ gives the following literal words in $D,V$:
\[
\begin{aligned}
 x&=D^3VDV^3,\\
 y&=D^3V^2DV^2,\\
 z_N&=D^2VDV^2DV.
\end{aligned}                                                    \tag{LI.7}
\]
The $27$ words $x^iy^jz_N^k$ are all different and exhaust
the kernel.  The commutator $z_N$ is central there.
The program multiplies every pair of those $27$ matrices
and verifies exactly (LI.3), for all $729$ products.
This proves both the nonabelian kernel and its exact law.
The explicit isomorphism to the group in (LI.2) is
\[
 (i,j,k)\longmapsto
 \begin{pmatrix}1&i&k+ij\\0&1&j\\0&0&1\end{pmatrix}.
\]
Multiplying these matrices gives (LI.3); their three upper
entries recover $i,j,k$ uniquely.

An actual complementary element is
\[
       v=V^2D^2VD^2V^2,\qquad u=D.                    \tag{LI.8}
\]
The sixteen $u^pv^q$ are distinct, commute, and have
exponents modulo four.  Multiplying all twenty-seven kernel
matrices by them gives the $432$ distinct matrices already
enumerated.  Direct conjugation of $x,y,z_N$ by $u,v$ gives
(LI.5).  This proves the exact semidirect-product description,
including the original generator coordinates (LI.4).
The quotient by $N$ is abelian, so $[Q,Q]\subseteq N$.
The action in (LI.5) on $N/Z(N)$ has matrix
$\left(\begin{smallmatrix}2&1\\1&1\end{smallmatrix}\right)$,
whose difference from the identity is invertible over
$\mathbb F_3$.  Consequently $[u,N]$ maps onto $N/Z(N)$.
Also $[x,y]=z_N$ lies in $[Q,Q]$.  These two facts prove
$N\subseteq[Q,Q]$, hence equality.

Both $u,v$ have the same action (LI.5), so $uv^{-1}$
is central of order four.  The central $z_N$ has order three.
Conversely, a central element $(n;u^pv^q)$ must have
$p+q=0\pmod4$, because the displayed action has order four
on $N/Z(N)$.  Centralizing $u$ then makes the image of $n$
in $N/Z(N)$ zero, since that action has no nonzero fixed
vector.  Thus the center is exactly
$\langle z_N\rangle\times\langle uv^{-1}\rangle\cong C_{12}$.
The complete centralizer enumeration independently gives
the same twelve matrices.  The coordinate
$(0,0,2;2,2)$ has order six, proving the assertion about
$(DV)^2$.  The orders of $D,V,DV$ are checked directly.
The normal-form multiplication gives (LI.6), thereby retaining
the precise two evaluation maps rather than only their target
isomorphism types.

The cusp interchange $D\leftrightarrow V$ descends to an
actual automorphism of $Q$.  The independent program checks it
on every vertex and every edge of the complete $432$-element
Cayley graph.  It sends
\[
 x\mapsto x^{-1}z_N,\qquad y\mapsto y^{-1}z_N,
                 \qquad z_N\mapsto z_N.                \tag{LI.10}
\]
Consequently it sends $(i,j,k)$ to $(-i,-j,k+i+j)$.
Applying it to $f(D,V^2)$ gives $f(V,D^2)$, which is the
inverse of $f(D^2,V)$ by antisymmetry.  This proves equality
of the first two coordinates of the two values.  Their
central coordinates also satisfy
$k_1+k_2=-i-j-ij$, directly from (LI.3).
Finally $D^4=V^4=q_{\rm loc}(a)^4=1$ and
$\omega a^2=(dV)^2$ reduce (ABR.4) exactly to
\[
                  f(D,V^2)=f(D^2,V)z^{-2c}.
                                                               \tag{LI.11}
\]
The coordinate of $z^{-2c}$ in (LI.4) is
$(0,0,2c;0,0)$.  This proves (LI.9), with no discarded
central coordinate.

The producer files are
\path{local_braid_quaternion_image.py},
\path{local_mixed_quaternion_census.py}, their generated GAP
programs and text outputs, and
\path{local_mixed_normal_form.py} with its JSON and NPZ
certificate.  The last program retains every mixed element,
all source words and the complete normal forms.  GAP's
structural identification package was unavailable in one
intermediate run; that failed request and its generated source
are retained separately.  The executed conclusions here use
explicit orders, commutators, powers, and the complete matrix
enumeration, and do not depend on that package.
\end{proof}

\paragraph{The exact limitation of the class-two ten-point target.}
At the class-two odd-exponent surface quotient, the pure
sphere mapping-class image lies in the abelian IA group
$\operatorname{Hom}(H_1,\bigwedge^2H_1)$ modulo its inner
subspace.  The label centralizer for the named ten-point
Klein-four configuration is $C_2^5$.  Its actual centralizer
image is consequently an extension of an abelian group by
an abelian group, and its derived subgroup is abelian.
A surjective homomorphism carries a derived subgroup onto
the derived subgroup of its target.  It therefore cannot
surject onto the full group $Q$ in (LI.2), whose derived
subgroup is the nonabelian group (LI.3).
This proves an obstruction to that particular full-target
factorization; it does not rule out a smaller conjugation
target or a deeper peripheral quotient.

For that distinction, the certificate also enumerates the
actual quotient $Q/Z(Q)$, of order thirty-six, with its two
generator maps, and the conjugacy class of $D$, of size nine.
Its centralizer has order forty-eight.  A defect in the full
(ABR.4) identity which is central in $Q$ disappears when
conjugating $D$; this is the exact quotient map relevant to
the weaker generator-image comparison.  The full weight
identity and this weaker conjugation condition are retained
as separate statements.  Equations (LI.9)--(LI.11) prove the
weaker conjugation condition from antisymmetry itself, because
the remaining defect is central.  Identifying this local
action with the action of a complete finite pentagon shadow
still requires the specified common source map; the full
central weight equation (LI.9) has not been removed.
